% Generated draft: requires mathematical and source review.
\documentclass[11pt]{article}
\usepackage[margin=25mm]{geometry}
\usepackage{fontspec}
\setmainfont{Latin Modern Roman}
\usepackage{amsmath,amssymb,mathtools,unicode-math}
\setmathfont{Latin Modern Math}
\usepackage{enumitem,needspace,xurl,hyperref}
\setlength{\parindent}{0pt}
\setlength{\parskip}{5pt}
\setlength{\emergencystretch}{3em}
\newcommand{\N}{\mathbb{N}}
\newcommand{\Z}{\mathbb{Z}}
\newcommand{\Q}{\mathbb{Q}}
\newcommand{\R}{\mathbb{R}}
\newcommand{\C}{\mathbb{C}}
\newcommand{\F}{\mathbb{F}}
\newcommand{\A}{\mathbb{A}}
\newcommand{\PP}{\mathbb{P}}
\newcommand{\E}{\mathbb{E}}
\newcommand{\eps}{\varepsilon}
\newcommand{\dd}{\,\mathrm{d}}
\begin{document}
\section*{UnsolvedMath ID: 700657}
\textbf{Source problem code:} catalogue entry
\subsection{Definitions and mathematical statement}
A group \(G\) is poly-\(\Z\) if it has a finite subnormal series
\[
\{1\}=G_0\trianglelefteq G_1\trianglelefteq\cdots\trianglelefteq G_r=G,
\qquad G_i/G_{i-1}\cong\Z\quad(1\leq i\leq r).
\]
The case \(r=0\) includes the trivial group. For a field \(k\), let \(R=kG\) be the group algebra: its elements are finite sums \(\sum_{g\in G}a_g g\), with coefficients in \(k\), and multiplication extends the group multiplication bilinearly.

Fix a module side, either left or right. An \(R\)-module \(P\) on that side is finitely generated projective if there exist an \(R\)-module \(Q\) on the same side and an integer \(n\geq0\) such that
\[
P\oplus Q\cong R^n.
\]
For left modules, a left ideal is an additive subgroup \(I\subseteq R\) satisfying \(RI\subseteq I\); for right modules, a right ideal satisfies \(IR\subseteq I\).

\textbf{Question (source formulation, with unresolved ideal convention).} Let \(G\) be any poly-\(\Z\) group and \(k\) any field. Is every finitely generated projective \(kG\)-module either free or isomorphic to an ideal?

If ``ideal'' means a one-sided ideal on the module's side, the corresponding precise question, for a fixed choice of side, is whether every such \(P\) satisfies
\[
\bigl(\exists m\in\Z_{\geq0}: P\cong R^m\bigr)
\quad\text{or}\quad
\bigl(\exists I\subseteq R:\ I\text{ is an ideal on that side and }P\cong I\bigr),
\]
where all isomorphisms are \(R\)-module isomorphisms on the chosen side. The supplied source specifies neither the module side nor whether ``ideal'' means a one-sided or a two-sided ideal. If a two-sided ideal is intended, the second alternative instead requires \(RI\subseteq I\) and \(IR\subseteq I\), with \(I\) viewed on the chosen module side. These conventions remain unresolved; the displayed one-sided interpretation is not an assertion that both side-specific questions were explicitly posed. The phrase ``is an ideal'' is interpreted here up to module isomorphism.
\subsection{Short English statement}
For a poly-\(\Z\) group \(G\) and any field \(k\), must every finitely generated projective module over \(kG\) be free or isomorphic to an ideal? The source leaves the module side and the meaning of ``ideal'' unspecified. The registry attributes the problem to G. A. Noskov, while the statement names J. T. Stafford.
\subsection{Sources}
\begin{itemize}
\item Local UnsolvedMath registry, supplied complete source record: permanent ID \(700657\), exact source problem code \texttt{catalogue entry}, title \emph{Kourovka Notebook Problem 10.42}; upstream revision \texttt{d371ee917ac9647fbb06cd6fc529deec6f3dda82}.
\item Supplied bibliographic reference: \emph{The Kourovka Notebook (Unsolved Problems in Group Theory)}, problem 10.42, 10th issue (1986), as identified in the supplied 21st-issue edition (2026), \texttt{arXiv:1401.0300v48}. The supplied collection description names E. I. Khukhro and V. D. Mazurov as editors. This reference has not been independently consulted here.
\item Attribution in the supplied background and proposer field: G. A. Noskov, proposed in 1986. The supplied statement separately begins with the attribution ``(J. T. Stafford)''; the relationship between these attributions is unspecified.
\item The supplied registry labels the problem \texttt{open} and records an update date of 2026-10-16. This is supplied registry metadata, not an independently verified assessment of the literature or mathematical status.
\end{itemize}
\end{document}
